%% kapa: exam=1 round=33 year=2022 period=1 form=A subject=real_estate subject_ko=부동산학원론 seq=4 number=84
%% source: Q-Net 감정평가사 기출문제 (cst00309) artlSeq=5209997 file=raw/1008/5209997/2022년도 제33회 감정평가사 1차시험 1교시(민법,경제학,부동산학).pdf
%% conversion: kordoc 4.24.51 -> tools/gen1.py (auto); review=proofread-sonnet; flags=none
\begin{question}{84}
부동산의 특성에 관한 설명으로 옳은 것의 개수는?
\begin{framedtext}
○ 용도의 다양성은 최유효이용을 선택할 수 있는 근거가 된다. \par
○ 인접성은 외부효과의 원인이 된다. \par
○ 분할·합병의 가능성은 부동산의 가치를 변화시킨다. \par
○ 부동성은 인근지역과 유사지역의 분류를 가능하게 한다. \par
○ 영속성은 부동산활동을 장기적으로 고려하게 한다.
\end{framedtext}
\begin{choices}
  \item 1
  \item 2
  \item 3
  \item 4
  \item 5
\end{choices}
\end{question}
\answer{5}
